\chapter{SemEval-2025 Task 10: Multilingual Characterization and Extraction of Narratives}
\label{sec:semeval25-task10}

\section{Overview}
SemEval-2025 Task 10 introduced a multilingual benchmark for the analysis of news narratives and entity framing across five languages and two high-impact domains (Ukraine–Russia War, Climate Change). The task is decomposed into entity framing, narrative classification, and narrative extraction subtasks, combining multi-label classification and evidence-grounded generation. Participating teams demonstrated the effectiveness of hierarchical prompting and targeted context selection; however, grounded explanation generation and label imbalance across languages remain unsolved challenges \cite{Piskorski2025, Fane2025, Singh2025, Fenu2025}.

We participated in this shared task (together with Diana Nurbakova) as competitors in Subtask 2 (Narrative Classification); the experiments reported here were carried out in that shared-task context.

This shared task asked teams to build systems that identify and extract dominant narratives from multilingual news, prioritising evidence-grounded explanations and robustness across languages. We ran multilingual evaluations to compare per-language behaviour.

\section{Task decomposition}
The shared task was split into three subtasks:
\begin{enumerate}
  \item \textbf{Entity Framing (Subtask 1)}: Given an article and named-entity mentions, assign main and fine-grained role labels to each mention (multi-label, multi-class role classification). Grounding to spans and role taxonomies was required. \cite{Piskorski2025}
  \item \textbf{Narrative Classification (Subtask 2)}: Assign articles one or more narrative labels from a two-level taxonomy (main narrative $\rightarrow$ subnarrative). This is a multilingual multi-label document classification problem. \cite{Piskorski2025}
  \item \textbf{Narrative Extraction (Subtask 3)}: For a given article and its dominant narrative, generate a concise (max.~80 words) free-text explanation supporting that narrative, grounded in evidence from the article. This subtask focused on faithfulness and coverage of extracted explanations. \cite{Piskorski2025}
\end{enumerate}

\section{Data and annotation}
Training and evaluation data were collected from online news sources spanning 2022-mid-2024. The dataset includes articles from a mix of mainstream and lower-credibility outlets; annotations follow multilingual guidelines and predefined taxonomies for entity roles and narratives. \cite{propaganda_site2025}

\section{Participation and headline results}
The task attracted broad participation; hundreds of teams registered and dozens submitted official test results. Representative high-performing approaches included:
\begin{itemize}
  \item \textbf{Fane et al. (``Fane'')} - zero-shot LLM prompting with a hierarchical decomposition for entity framing; reported Main Role Accuracy $\approx$ 89.4\% and Exact Match Ratio $\approx$ 34.5\% on English. \cite{Fane2025}
  \item \textbf{GateNLP} - \emph{Hierarchical Three-Step Prompting (H3Prompt)} for Subtask 2 (narrative classification): cascade of domain $\rightarrow$ main narrative $\rightarrow$ subnarrative prompting; top performer on the English test set among competing teams. Code published. \cite{Singh2025}
  \item \textbf{DEMON / TechSSN / LTG / BERTastic} - representative competitive systems using fine-tuning of large pre-trained encoders (e.g., XLM-R / DeBERTa variants) or fine-tuned LLaMA-family models for framing and narrative tasks. Several system descriptions and official proceedings papers document model design choices and per-language performance differences. \cite{Fenu2025, Premnath2025, Ronningstad2025, Nakov2025}
\end{itemize}
These system papers illustrate a common trend: hierarchical or multi-step prompting and careful context selection (to fit model context windows) were effective strategies across teams. \cite{Fane2025, Singh2025, Ronningstad2025}

\section{Evaluation of Subtask 2: Narrative Classification}
\label{subsec:evaluation-subtask2}

The second subtask of SemEval-2025 Task 10 required participants to assign
\emph{narrative} (coarse-level) and \emph{sub-narrative} (fine-level) labels
to each news article. Since articles can be associated with multiple narratives,
this constitutes a multi-label, multi-class document classification problem
\cite{Piskorski2025}.

\paragraph{Evaluation target.}
Participants had to provide two sets of predictions per article:
(1) a set of coarse-level narrative labels and
(2) a set of fine-grained sub-narrative labels.
Although both levels were evaluated, the official ranking was based on performance at the sub-narrative level, as this captures more nuanced framing
\cite{propaganda_site2025}.

\paragraph{Metrics.}
Two main evaluation metrics were applied:
\begin{itemize}
  \item \textbf{Macro F1 (coarse level).}
  This is the unweighted mean of F1-scores computed for each coarse narrative class. It penalizes models that ignore low-frequency labels.
  \item \textbf{F1 (samples).}
  This metric is computed at the level of individual articles. For each document, precision, recall, and F1 are compute by comparing the predicted label set with the gold label set. The final score is the average of these values across all documents.
\end{itemize}
In addition, standard deviations across multiple runs and across languages were reported on the leaderboard to quantify robustness \cite{leaderboard2025}.

\paragraph{Worked example.}
Consider a simplified case with three possible sub-narratives:
$\{N_{1}, N_{2}, N_{3}\}$.
Suppose the gold labels for a document are $\{N_{1}, N_{3}\}$,
while a system predicts $\{N_{1}, N_{2}\}$.

\begin{itemize}
  \item \textbf{Precision} $= \frac{|\{N_{1}, N_{2}\} \cap \{N_{1}, N_{3}\}|}
  {|\{N_{1}, N_{2}\}|} = \frac{1}{2} = 0.5$.
  \item \textbf{Recall} $= \frac{|\{N_{1}, N_{2}\} \cap \{N_{1}, N_{3}\}|}
  {|\{N_{1}, N_{3}\}|} = \frac{1}{2} = 0.5$.
  \item \textbf{F1 (sample)} $= \frac{2 \times 0.5 \times 0.5}{0.5 + 0.5} = 0.5$.
\end{itemize}

In contrast, if the system had predicted only $\{N_{1}\}$,
precision would be $1.0$, recall $0.5$, and F1 (sample) $0.67$.
This illustrates that partial but precise predictions may yield a higher F1
than predicting extra incorrect labels.

\paragraph{Implications.}
Because the leaderboard ranks systems by fine-level labels using sample-based
F1, models must balance recall (capturing all relevant sub-narratives) with
precision (avoiding irrelevant ones). Macro F1 further incentivizes attention
to rare narratives, since underperformance on them disproportionately lowers
the score. \cite{Piskorski2025}.
